\documentclass[12pt]{article}
\usepackage[T1]{fontenc}
\usepackage[utf8]{inputenc}
\usepackage{lmodern}
\usepackage{textcomp}
\usepackage{enumitem}
\usepackage{geometry}
\usepackage{graphicx}
\usepackage{amsmath, amssymb}
\geometry{margin=1in}
\begin{document}
\begin{center}
Philosophy - Graduate Level Assessment 7 \
Total Marks: 40 \
Answer all questions.
\end{center}

\section*{Multiple Choice (10 marks)}
\begin{enumerate}[label=\arabic*.]
\item According to Parfit, the obligation to give priority to the welfare of one's children is:
\begin{enumerate}[label=(\alph*)]
    \item agent-relative.
    \item agent-neutral.
    \item absolute.
    \item none of the above.
\end{enumerate}
\item Mill claims that one of the strongest objections to utilitarianism is drawn from the idea of:
\begin{enumerate}[label=(\alph*)]
    \item justice.
    \item supererogation.
    \item honesty.
    \item morality.
    \item virtue.
\end{enumerate}
\item Sacred literature originated with which of the following jina?
\begin{enumerate}[label=(\alph*)]
    \item Sri Lakshmi
    \item Brahma
    \item Nanak
    \item Vishnu
    \item Shiva
\end{enumerate}
\item At first Descartes supposes that everything he sees is \_\_\_\_\_.
\begin{enumerate}[label=(\alph*)]
    \item TRUE
    \item undeniable
    \item FALSE
    \item unchangeable
    \item an illusion
\end{enumerate}
\item Nussbaum claims that in cross-cultural communication, inhabitants of different conceptual schemes tend to view their interaction:
\begin{enumerate}[label=(\alph*)]
    \item in a utilitarian way.
    \item in a Cartesian way.
    \item in a nihilistic way.
    \item in a Hegelian way.
    \item in a Kantian way.
\end{enumerate}
\end{enumerate}

\section*{True / False (10 marks)}
\textit{Answer True or False}
\begin{enumerate}[label=\arabic*.]
\setcounter{enumi}{5}
\item True or False: According to Hume, morality is ultimately based on cultural norms shaped by societal interactions and experiences.
\item True or False: Arthur contends that Singer overlooks the principles of equality and harm when discussing world hunger and poverty.
\item True or False: Pence argues that those who object to SCNT assume that parents are motivated by genuine concern for their child's wellbeing.
\item True or False: Papadaki actually agrees with Kant's view on the moral implications of sexual objectification.
\item True or False: Caityavasis refer to lay followers who primarily reside in urban centers and support monastic communities.
\end{enumerate}

\section*{Long Form Response (20 marks)}
\textit{Answer each question with a clear and complete explanation.}
\begin{enumerate}[label=\arabic*.]
\setcounter{enumi}{10}
\item Discuss the concept of intrinsic goods in Aquinas's natural law theory, analyzing their significance and interrelation within his ethical framework.
\item Discuss Feinberg's perspective on human desires in the context of hunger, analyzing why he posits that these desires do not align with traditional classifications.
\end{enumerate}

\vfill
\noindent\textit{End of Paper}
\end{document}